%% ============================================================
%% Teil III: Echtzeit und Performance (Kapitel 9–11)
%% Band 3 – Kapitel 9 bis 11
%% ============================================================

\part{Echtzeit und Performance}

% ─────────────────────────────────────────────────────────────
\chapter{Wir wollen wissen, wann C++ lohnt –
  Python vs. C++ für Robotik-Nodes}
% ─────────────────────────────────────────────────────────────

\begin{lernziele}
Du kannst begründen, welcher Node in Python und welcher in C++ sinnvoller ist,
kennst die tatsächlichen Latenz-Unterschiede
und verstehst das Hybridmodell von robotik-uno-q.
\end{lernziele}

\section{Situation}

Manche behaupten: "`Alles in C++ -- wegen Performance."'
Andere: "`Alles in Python -- wegen Entwicklungsgeschwindigkeit."'
Beide liegen halb falsch.

\begin{aufgabeBox}
Entscheide für \texttt{robotik-uno-q}, welcher Node
welche Sprache braucht -- mit konkreten Messungen.
\end{aufgabeBox}

\section{Was Python langsamer macht}

\begin{wissenBox}
\textbf{Der Global Interpreter Lock (GIL):}
Python kann nur in einem Betriebssystem-Thread gleichzeitig
Python-Bytecode ausführen.
Für IO-bound-Tasks (Serial lesen, warten) ist das kein Problem.
Für CPU-bound-Tasks (1000-Hz-Regler, Bildverarbeitung) schon.

\textbf{Garbage Collector:}
Pythons Speicherverwaltung kann unvorhersehbare Pausen
(GC-Pausen) verursachen -- typisch 1--50\,ms.
Für Echtzeit-Regler ist das inakzeptabel.

\textbf{Interpreter-Overhead:}
Jede Python-Operation löst mehrere C-Aufrufe aus.
Für \texttt{arduino\_bridge\_node} mit 10\,Hz ist das irrelevant.
Für einen 1\,kHz-IMU-Regler würde es auffallen.
\end{wissenBox}

\section{Latenz-Messung: Python vs. C++}

\begin{lstlisting}[language=bash]
# Timer-Jitter messen: wie präzise ist ein 100-Hz-Timer?
# Python (rclpy):
ros2 run robotik_uno_q arduino_bridge_node
ros2 topic hz /sensor/rechts --window 100
# average rate: 9.991  -- Jitter: ~1 ms

# C++ (rclcpp):
ros2 run robotik_uno_q_cpp motor_node
ros2 topic hz /motor/status --window 100
# average rate: 9.999  -- Jitter: <0.1 ms
\end{lstlisting}

Bei 10\,Hz ist der Unterschied für \texttt{robotik-uno-q} nicht spürbar.
Ab 100\,Hz wird C++ relevant.

\section{Entscheidungsbaum}

\begin{center}
\begin{tikzpicture}[
    node distance=1.8cm,
    auto,
    decision/.style={diamond, draw=orange!80!black, fill=orange!10,
                     minimum width=3cm, minimum height=1cm,
                     aspect=2, font=\small},
    result/.style={rectangle, draw=green!60!black, fill=green!10,
                   minimum width=3cm, font=\small}
]
\node[decision] (rate) {Rate > 100\,Hz?};
\node[decision, below left=of rate, xshift=-1cm] (rt) {Echtzeit-Anforderung?};
\node[decision, below right=of rate, xshift=1cm] (cpu) {CPU-intensiv?};
\node[result, below of=rt, yshift=-0.5cm] (cpp1) {C++};
\node[result, below of=cpu, yshift=-0.5cm] (py) {Python};
\node[result, below right=of rt, xshift=0.5cm, yshift=-0.5cm] (cpp2) {C++};

\draw[->] (rate) -- node[left]{Ja} (rt);
\draw[->] (rate) -- node[right]{Nein} (cpu);
\draw[->] (rt)   -- node[left]{Ja}  (cpp1);
\draw[->] (rt)   -- node[right]{Nein} (cpp2);
\draw[->] (cpu)  -- node[right]{Nein} (py);
\draw[->] (cpu)  -- node[left]{Ja}  (cpp2);
\end{tikzpicture}
\end{center}

\textbf{Für robotik-uno-q:}

\begin{center}
\begin{tabular}{lll}
\hline
\textbf{Node} & \textbf{Rate} & \textbf{Sprache} \\
\hline
\texttt{arduino\_bridge\_node} & 10\,Hz Serial & Python \\
\texttt{navigation\_node}      & 10\,Hz Entscheidung & Python \\
\texttt{motor\_node}           & 10\,Hz + Service & C++ \\
\hline
\end{tabular}
\end{center}

\texttt{motor\_node} in C++ nicht wegen Performance,
sondern wegen Lifecycle~Node -- das gibt es in rclcpp vollständiger.

\begin{projektBox}[robotik-uno-q -- Projektstand]
\textbf{Heute:} Begründete Sprachentscheidung für jeden Node.\\
\textbf{Nächstes Kapitel:} Echtzeit-Konfiguration des rclcpp-Executors.
\end{projektBox}

\begin{merksatzBox}
\textbf{Merksatz:}
C++ für Rates über 100\,Hz, Echtzeit-Anforderungen und CPU-intensive Berechnungen.
Python für alles andere -- und das ist in der Praxis die Mehrheit der Nodes.
Das Hybridmodell ist kein Kompromiss, sondern die korrekte Architektur.
\end{merksatzBox}


% ─────────────────────────────────────────────────────────────
\chapter{Wir wollen vorhersagbare Laufzeiten –
  Echtzeit mit rclcpp}
% ─────────────────────────────────────────────────────────────

\begin{lernziele}
Du verstehst, was Echtzeit bedeutet (nicht: schnell, sondern: vorhersagbar),
konfigurierst den rclcpp-Executor für niedrigen Jitter
und kennst die PREEMPT\_RT-Verbindung aus Band~1.
\end{lernziele}

\section{Situation}

Ein Node published Steuerkommandos an 100\,Hz.
Manchmal kommt ein Tick nach 10\,ms, manchmal nach 15\,ms, manchmal nach 8\,ms.
Das ist kein Echtzeit-System.
Ein Echtzeit-System liefert jeden Tick in 10\,ms~$\pm$\,0{,}1\,ms.

\begin{aufgabeBox}
Konfiguriere den \texttt{motor\_node} für minimalen Jitter:
\texttt{StaticSingleThreadedExecutor}, Speicher-Präallokation,
und PREEMPT\_RT-Scheduling aus Band~1.
\end{aufgabeBox}

\section{Executor-Typen in rclcpp}

\begin{wissenBox}
\textbf{SingleThreadedExecutor} (Standard):
Alle Callbacks in einem Thread, sequenziell.
Einfach, aber: ein langer Callback blockiert alle anderen.

\textbf{MultiThreadedExecutor}:
Callbacks in mehreren Threads parallel.
Höherer Durchsatz, aber: Callbacks müssen thread-safe sein.

\textbf{StaticSingleThreadedExecutor} (Echtzeit):
Ähnlich Single, aber Callback-Gruppen werden einmalig beim Start
alloziert -- keine dynamischen Allokationen zur Laufzeit.
Das verhindert GC-artige Pausen.

\begin{lstlisting}[style=cpp]
int main(int argc, char* argv[]) {
    rclcpp::init(argc, argv);
    auto node = std::make_shared<MotorNode>();

    // StaticSingleThreadedExecutor: keine Runtime-Allokationen
    rclcpp::executors::StaticSingleThreadedExecutor executor;
    executor.add_node(node->get_node_base_interface());
    executor.spin();

    rclcpp::shutdown();
    return 0;
}
\end{lstlisting}
\end{wissenBox}

\section{Speicher-Präallokation}

\begin{lstlisting}[style=cpp]
// Im Konstruktor: Puffer vorab allozieren
MotorNode() : LifecycleNode("motor") {
    // Vector mit fester Kapazität vorab allozieren
    // Keine Runtime-Allokation beim ersten Callback
    puffer_.reserve(100);
    puffer_.resize(100, 0.0f);
}
private:
    std::vector<float> puffer_;  // vorab alloziert
\end{lstlisting}

\begin{lstlisting}[language=bash]
# PREEMPT_RT-Scheduling aus Band 1 aktivieren
# (erfordert PREEMPT_RT-Kernel, in Band 1 Kapitel 29 eingerichtet)
sudo chrt -f 80 ros2 run robotik_uno_q_cpp motor_node
# -f: FIFO-Scheduling, Priorität 80 (höher = höher priorisiert)
\end{lstlisting}

\begin{projektBox}[robotik-uno-q -- Projektstand]
\textbf{Heute:} \texttt{motor\_node} mit \texttt{StaticSingleThreadedExecutor}.\\
\textbf{Nächstes Kapitel:} Profiling -- Flaschenhälse finden.
\end{projektBox}

\begin{merksatzBox}
\textbf{Merksatz:}
Echtzeit $\neq$ schnell. Echtzeit = vorhersagbar.
\texttt{StaticSingleThreadedExecutor} verhindert Runtime-Allokationen.
Speicher vorab allozieren (\texttt{reserve()}).
PREEMPT\_RT-Kernel + \texttt{chrt} für echtes Echtzeit-Scheduling.
\end{merksatzBox}


% ─────────────────────────────────────────────────────────────
\chapter{Wir wollen Flaschenhälse finden –
  Profiling und Optimierung}
% ─────────────────────────────────────────────────────────────

\begin{lernziele}
Du verwendest \texttt{ros2 topic hz}, \texttt{perf} und
rclcpp-interne Latenz-Messung, um den Node-Stack zu profilieren
und Engpässe zu finden.
\end{lernziele}

\section{Situation}

\texttt{ros2 topic hz /cmd\_vel} zeigt 7\,Hz statt 10\,Hz.
Irgendetwas verlangsamt den Stack.
Profiling findet, was.

\begin{aufgabeBox}
Profiliere den \texttt{motor\_node} und finde, ob der Engpass
im Serial-Schreiben, im Twist-Parsing oder anderswo liegt.
\end{aufgabeBox}

\section{Werkzeuge}

\subsection{ros2 topic hz und delay}

\begin{lstlisting}[language=bash]
# Frequenz messen
ros2 topic hz /cmd_vel --window 50
# average rate: 9.87  min: 0.099  max: 0.104  std dev: 0.001

# Latenz: Zeitstempel in Nachricht vs. Empfangszeitpunkt
ros2 topic delay /cmd_vel
# average delay: 0.003  -- 3 ms Latenz, akzeptabel
\end{lstlisting}

\subsection{Manuelle Latenz-Messung im Node}

\begin{lstlisting}[style=cpp]
void _verarbeite_twist(const Twist::SharedPtr& msg) {
    auto t0 = std::chrono::high_resolution_clock::now();

    _sende_serial(/* ... */);

    auto t1 = std::chrono::high_resolution_clock::now();
    auto us = std::chrono::duration_cast<std::chrono::microseconds>(t1 - t0).count();

    if (us > 1000) {
        RCLCPP_WARN(get_logger(), "Serial-Latenz: %ld us -- zu hoch!", us);
    }
}
\end{lstlisting}

\subsection{perf für System-Level-Profiling}

\begin{lstlisting}[language=bash]
# perf: CPU-Zyklen pro Funktion messen
# (nur auf Systemen mit Perf-Support -- nicht im Container ohne --privileged)
perf record -g ros2 run robotik_uno_q_cpp motor_node
perf report
# Zeigt: welche Funktion wie viel CPU-Zeit verbraucht
\end{lstlisting}

\section{Häufige Engpässe und Lösungen}

\begin{center}
\begin{tabular}{ll}
\hline
\textbf{Engpass} & \textbf{Lösung} \\
\hline
Serial-Timeout zu hoch & \texttt{timeout=0.1} statt \texttt{1.0} \\
\texttt{dynamic\_cast} im Callback & \texttt{static\_cast} oder Template \\
String-Formatierung im Callback & Buffer vorallozieren, \texttt{snprintf} \\
\texttt{shared\_ptr}-Kopien & \texttt{const SharedPtr\&} statt \texttt{SharedPtr} \\
Python-Node als Flaschenhals & zu C++ migrieren, wenn Jitter > 5\,ms \\
\hline
\end{tabular}
\end{center}

\begin{projektBox}[robotik-uno-q -- Projektstand]
\textbf{Heute:} Profiling-Werkzeuge eingesetzt, Serial-Latenz unter 500\,µs.\\
\textbf{Teil~IV:} Das vollständige Projekt -- C++ und Python im gemeinsamen Stack.
\end{projektBox}

\begin{merksatzBox}
\textbf{Merksatz:}
Immer erst messen, dann optimieren.
\texttt{ros2 topic hz} und \texttt{delay} sind die ersten Werkzeuge.
Manuelle \texttt{chrono}-Messung im Callback findet den Engpass im Node.
\texttt{perf} für System-Level -- braucht Root-Rechte und echte Hardware.
\end{merksatzBox}
